\documentclass[12pt,a4paper]{article}

\usepackage[utf8]{inputenc}
\usepackage[T1]{fontenc}
\usepackage[french]{babel}
\usepackage{geometry}
\geometry{top=2.5cm, bottom=2.5cm, left=2.5cm, right=2.5cm, headheight=15pt}
\usepackage{booktabs}
\usepackage{array}
\usepackage{tabularx}
\usepackage{longtable}
\usepackage{xcolor}
\usepackage{fancyhdr}
\usepackage{setspace}
\usepackage{enumitem}
\usepackage{hyperref}
\usepackage{microtype}

\hypersetup{
    colorlinks=true,
    linkcolor=black,
    urlcolor=blue,
    pdftitle={Journal de Stage -- CHRONO\_01},
    pdfauthor={Becem Ameur},
}

\pagestyle{fancy}
\fancyhf{}
\fancyhead[L]{Journal de Stage -- Becem Ameur}
\fancyhead[R]{ENIM -- 2025/2026}
\fancyfoot[C]{\thepage}
\renewcommand{\headrulewidth}{0.4pt}
\renewcommand{\footrulewidth}{0.4pt}

\onehalfspacing

\definecolor{semaine1}{rgb}{0.85,0.93,1.0}
\definecolor{semaine2}{rgb}{0.85,1.0,0.88}
\definecolor{semaine3}{rgb}{1.0,0.96,0.82}
\definecolor{semaine4}{rgb}{0.96,0.85,1.0}
\definecolor{semaine5}{rgb}{1.0,0.88,0.85}
\definecolor{titlebg}{rgb}{0.18,0.36,0.60}

% ==============================================================================
\begin{document}
% ==============================================================================

% ----- Titre ------------------------------------------------------------------
\begin{titlepage}
    \begin{center}
        \vspace*{1.5cm}
        {\Large \textbf{École Nationale d'Ingénieurs de Monastir (ENIM)}}\\[0.3cm]
        {\large Université de Monastir -- Tunisie}\\[1.5cm]
        \rule{\linewidth}{1.2pt}\\[0.4cm]
        {\huge \textbf{Journal de Stage}}\\[0.3cm]
        \rule{\linewidth}{1.2pt}\\[1.0cm]
        {\LARGE \textbf{Projet CHRONO\_01}}\\[0.4cm]
        {\large Chronomètre géant à afficheurs 7 segments 24\,V\\
        avec contrôle WiFi et GSM}\\[2.0cm]

        \begin{tabular}{ll}
            \textbf{Stagiaire :}    & Becem Ameur \\[0.2cm]
            \textbf{Spécialité :}   & Génie Électrique, 3\textsuperscript{ème} année \\[0.2cm]
            \textbf{Établissement :}& ENIM, Université de Monastir \\[0.2cm]
            \textbf{Durée :}        & 30 jours (6 semaines) \\[0.2cm]
            \textbf{Année :}        & 2025 -- 2026 \\
        \end{tabular}

        \vfill
        \rule{\linewidth}{0.6pt}\\[0.2cm]
        {\small Document réservé à l'usage académique -- ENIM 2026}
    \end{center}
\end{titlepage}

\newpage
\tableofcontents
\newpage

% ==============================================================================
\section*{Introduction}
\addcontentsline{toc}{section}{Introduction}
% ==============================================================================

Ce journal retrace toutes les activités que j'ai réalisées pendant mon stage de 30 jours sur le projet \textbf{CHRONO\_01}. L'objectif de ce projet était de concevoir et de fabriquer un chronomètre géant à afficheurs 7 segments alimentés en 24\,V, avec un système de contrôle sans fil (WiFi et GSM).

J'ai organisé mon travail en cinq grandes semaines, en commençant par la recherche et la planification, puis en passant à la conception électronique, à la réalisation du circuit imprimé (PCB), à la programmation, et enfin aux tests et à la mise en boîtier.

\bigskip
\noindent\textbf{Découpage général du stage :}
\begin{itemize}
    \item \textbf{Semaine 1 (Jours 1--6) :} Recherche, analyse du besoin et choix des composants
    \item \textbf{Semaine 2 (Jours 7--12) :} Conception du schéma électrique sous KiCad 10
    \item \textbf{Semaine 3 (Jours 13--18) :} Routage du PCB et préparation des fichiers de fabrication
    \item \textbf{Semaine 4 (Jours 19--24) :} Programmation du firmware et interface web
    \item \textbf{Semaine 5--6 (Jours 25--30) :} Assemblage mécanique, tests finaux et documentation
\end{itemize}

\newpage

% ==============================================================================
\section{Semaine 1 -- Recherche et planification (Jours 1 à 6)}
% ==============================================================================

\colorbox{semaine1}{\parbox{\linewidth}{\centering\textbf{Objectif : Comprendre le projet, choisir les composants et préparer le cahier des charges}}}

\bigskip

\subsection*{Jour 1 -- Présentation et analyse du projet}
\begin{itemize}
    \item Lecture des informations disponibles sur les afficheurs 7 segments de grande taille.
    \item Dessin d'un premier schéma de principe à la main : 6 chiffres, chacun avec 7 segments = 42 LED.
    \item Définition des besoins principaux : alimentation 24\,V, contrôle sans fil, boîtier solide.
    \item Rédaction d'une première liste de questions techniques à résoudre.
\end{itemize}

\subsection*{Jour 2 -- Recherche sur les microcontrôleurs}
\begin{itemize}
    \item Étude comparative entre plusieurs microcontrôleurs : Arduino, STM32, ESP32.
    \item Choix de l'\textbf{ESP32-WROVER-E} car il a le WiFi et le Bluetooth intégrés, assez de mémoire et beaucoup de broches GPIO.
    \item Lecture du datasheet de l'ESP32-WROVER-E pour comprendre ses caractéristiques.
    \item Comparaison avec d'autres options et justification du choix final.
\end{itemize}

\subsection*{Jour 3 -- Recherche sur la commande des segments LED}
\begin{itemize}
    \item Problème identifié : l'ESP32 n'a pas assez de sorties pour commander 42 segments directement.
    \item Solution trouvée : utiliser des \textbf{registres à décalage 74HC595} en cascade (daisy-chain).
    \item Calcul du nombre de registres nécessaires : 42 segments $\div$ 8 sorties/registre = 6 registres.
    \item Compréhension du protocole SPI utilisé pour envoyer les données aux registres.
\end{itemize}

\subsection*{Jour 4 -- Choix des circuits d'amplification de courant}
\begin{itemize}
    \item Les sorties du 74HC595 ne peuvent pas fournir assez de courant pour les LED 24\,V.
    \item Recherche et sélection des \textbf{ULN2803A} : tableaux de transistors Darlington qui peuvent soutirer jusqu'à 500\,mA par canal.
    \item Lecture de la fiche technique de l'ULN2803A.
    \item Calcul du courant total maximal du système : 42 segments × 150\,mA = 6,3\,A.
\end{itemize}

\subsection*{Jour 5 -- Alimentation et module GSM}
\begin{itemize}
    \item Choix de l'alimentation principale : transformateur 230\,V / 24\,V avec filtre redresseur.
    \item Sélection du convertisseur abaisseur \textbf{LM2596} pour passer de 24\,V à 5\,V (pour l'ESP32 et les circuits logiques).
    \item Étude du module \textbf{SIM800L} pour la communication GSM (envoi/réception de SMS).
    \item Identification d'un problème critique : le SIM800L consomme jusqu'à 2\,A en pointe lors des transmissions.
\end{itemize}

\subsection*{Jour 6 -- Calculs électriques et rédaction du cahier des charges}
\begin{itemize}
    \item Calcul du condensateur de filtrage pour le SIM800L : $C_1 = 1000\,\mu\text{F}$.
    \item Calcul du condensateur d'entrée du LM2596 : $C_2 = 100\,\mu\text{F}$.
    \item Rédaction du tableau des spécifications techniques du projet (cahier des charges).
    \item Préparation d'une liste définitive des composants à commander (BOM).
\end{itemize}

\newpage

% ==============================================================================
\section{Semaine 2 -- Conception du schéma électrique (Jours 7 à 12)}
% ==============================================================================

\colorbox{semaine2}{\parbox{\linewidth}{\centering\textbf{Objectif : Dessiner le schéma électrique complet sous KiCad 10}}}

\bigskip

\subsection*{Jour 7 -- Prise en main de KiCad 10}
\begin{itemize}
    \item Installation et découverte du logiciel KiCad 10.
    \item Création d'un nouveau projet : \texttt{CHRONO\_01}.
    \item Apprentissage de l'interface de l'éditeur de schémas (Schematic Editor).
    \item Placement des premiers symboles : résistances, condensateurs, connecteurs.
\end{itemize}

\subsection*{Jour 8 -- Placement des composants principaux}
\begin{itemize}
    \item Ajout du symbole de l'ESP32-WROVER-E dans le schéma.
    \item Placement des 6 registres à décalage 74HC595 et connexion en cascade SPI.
    \item Placement des 6 réseaux de transistors ULN2803A.
    \item Connexion des signaux SPI : \texttt{CLK}, \texttt{DATA}, \texttt{LATCH}.
\end{itemize}

\subsection*{Jour 9 -- Schéma d'alimentation}
\begin{itemize}
    \item Dessin du circuit d'alimentation : entrée 24\,V, convertisseur LM2596, sortie 5\,V.
    \item Ajout des condensateurs de découplage sur chaque circuit intégré.
    \item Placement du condensateur de tampon $C_1 = 1000\,\mu\text{F}$ pour le SIM800L.
    \item Ajout de la LED de puissance $D_1$ avec sa résistance de limitation $R_1 = 2{,}4\,\text{k}\Omega$.
\end{itemize}

\subsection*{Jour 10 -- Module GSM et connecteurs}
\begin{itemize}
    \item Ajout du module SIM800L dans le schéma avec ses connexions UART (TX, RX).
    \item Placement des connecteurs de sortie pour les 6 chiffres (connecteurs JST-XH 8 broches).
    \item Ajout du connecteur d'entrée principal d'alimentation 24\,V (bornier à vis).
    \item Création des points de test (TP1, TP2) pour les mesures oscilloscope.
\end{itemize}

\subsection*{Jour 11 -- Vérification du schéma (ERC)}
\begin{itemize}
    \item Lancement de la vérification électrique automatique (ERC -- Electrical Rules Check).
    \item Correction de toutes les erreurs trouvées : fils non connectés, broches manquantes.
    \item Ajout des drapeaux \texttt{PWR\_FLAG} sur les rails d'alimentation.
    \item Vérification que chaque broche de chaque CI est bien connectée ou terminée correctement.
\end{itemize}

\subsection*{Jour 12 -- Finalisation du schéma et empreintes}
\begin{itemize}
    \item Assignation des empreintes (footprints) à chaque composant du schéma.
    \item Choix des empreintes DIP-16 pour les 74HC595 et DIP-18 pour les ULN2803A.
    \item Export de la liste des connexions (netlist) vers le routeur PCB.
    \item Sauvegarde et archivage du projet KiCad.
\end{itemize}

\newpage

% ==============================================================================
\section{Semaine 3 -- Routage du PCB et fabrication (Jours 13 à 18)}
% ==============================================================================

\colorbox{semaine3}{\parbox{\linewidth}{\centering\textbf{Objectif : Dessiner le circuit imprimé (PCB), le faire fabriquer et le souder}}}

\bigskip

\subsection*{Jour 13 -- Découverte de l'éditeur PCB de KiCad}
\begin{itemize}
    \item Ouverture de l'éditeur PCB (PCB Editor) dans KiCad 10.
    \item Import des composants depuis le schéma électrique.
    \item Définition des dimensions du PCB : premiers essais avec un format 280 × 110\,mm.
    \item Étude des contraintes de fabrication artisanale (trous non métallisés, NPTH).
\end{itemize}

\subsection*{Jour 14 -- Placement des composants sur le PCB}
\begin{itemize}
    \item Organisation des composants par blocs logiques : un bloc par chiffre.
    \item Mise au point d'une géométrie modulaire en forme de ``T'' pour chaque chiffre.
    \item Réduction de la taille totale du PCB à 160 × 80\,mm (gain de 58\,\%).
    \item Placement des sockets tulipes à broches usinées pour les circuits DIP.
\end{itemize}

\subsection*{Jour 15 -- Routage des pistes}
\begin{itemize}
    \item Routage des pistes d'alimentation (24\,V) avec une largeur de 2\,mm.
    \item Routage des pistes SPI et de commande avec une largeur de 0,6\,mm.
    \item Respect des distances minimales entre pistes : $\geq 0{,}38$\,mm.
    \item Ajout des zones de cuivre (copper pours) sur les deux faces du PCB pour la masse.
\end{itemize}

\subsection*{Jour 16 -- Vérification DRC et optimisation}
\begin{itemize}
    \item Lancement de la vérification des règles de conception (DRC -- Design Rules Check).
    \item Correction de toutes les violations trouvées (0 erreur obtenue à la fin).
    \item Ajout des gouttes (teardrops) sur les pastilles pour solidifier les soudures.
    \item Vérification visuelle de la vue 3D du PCB.
\end{itemize}

\subsection*{Jour 17 -- Génération des fichiers et envoi à la fabrication}
\begin{itemize}
    \item Génération des fichiers Gerber (RS-274X) pour chaque couche du PCB.
    \item Génération des fichiers de perçage Excellon (.drl).
    \item Vérification des fichiers avec un viewer Gerber.
    \item Envoi des fichiers à l'artisan local pour la gravure chimique. Coût : 120\,TND.
\end{itemize}

\subsection*{Jour 18 -- Réception du PCB et soudure des composants}
\begin{itemize}
    \item Réception du PCB fabriqué et vérification visuelle (pas de court-circuit visible).
    \item Soudure des composants passifs : résistances, condensateurs.
    \item Installation des sockets tulipes pour les circuits intégrés DIP.
    \item Soudure des borniers et connecteurs JST.
\end{itemize}

\newpage

% ==============================================================================
\section{Semaine 4 -- Programmation du firmware (Jours 19 à 24)}
% ==============================================================================

\colorbox{semaine4}{\parbox{\linewidth}{\centering\textbf{Objectif : Écrire le programme de l'ESP32 et créer l'interface web}}}

\bigskip

\subsection*{Jour 19 -- Mise en place de l'environnement de développement}
\begin{itemize}
    \item Installation de l'IDE Arduino et de la bibliothèque ESP32 (ESP32 Arduino Core).
    \item Test de connexion entre le PC et la carte ESP32-WROVER-E.
    \item Premier programme de test : allumer et éteindre tous les segments (test de câblage).
    \item Vérification que les 42 segments répondent correctement aux commandes.
\end{itemize}

\subsection*{Jour 20 -- Programmation des registres à décalage (SPI)}
\begin{itemize}
    \item Écriture de la fonction qui envoie les données aux 6 registres 74HC595 en cascade via SPI.
    \item Création d'un tableau de correspondance (look-up table) pour les chiffres 0 à 9.
    \item Test d'affichage de chaque chiffre sur chacun des 6 digit du chronomètre.
    \item Vérification visuelle que les bonnes LED s'allument pour chaque chiffre.
\end{itemize}

\subsection*{Jour 21 -- Logique du chronomètre}
\begin{itemize}
    \item Écriture de la logique de comptage : heures, minutes, secondes.
    \item Gestion des dépassements (59 secondes $\rightarrow$ 0 + 1 minute, etc.).
    \item Utilisation d'une tâche sur le cœur 1 de l'ESP32 (dual-core) pour le comptage.
    \item Tests de précision du temps avec un chronomètre de référence.
\end{itemize}

\subsection*{Jour 22 -- Serveur web WiFi (interface de contrôle)}
\begin{itemize}
    \item Configuration de l'ESP32 en mode Point d'Accès WiFi (Soft-AP).
    \item Création d'un serveur HTTP asynchrone avec la bibliothèque \texttt{ESPAsyncWebServer}.
    \item Écriture de la page web HTML/CSS/JavaScript embarquée dans le firmware.
    \item Test de la page web depuis un téléphone connecté au réseau WiFi du chronomètre.
\end{itemize}

\subsection*{Jour 23 -- Commandes web et ajustement de l'heure}
\begin{itemize}
    \item Ajout des boutons sur la page web : Démarrer, Arrêter, Remettre à zéro.
    \item Ajout d'un formulaire pour régler l'heure manuellement (HH:MM:SS).
    \item Test de chaque bouton depuis l'interface web sur téléphone et tablette.
    \item Correction d'un problème d'affichage : les deux-points (:) ne clignotaient pas correctement.
\end{itemize}

\subsection*{Jour 24 -- Module GSM (SIM800L) et commande par SMS}
\begin{itemize}
    \item Câblage du module SIM800L sur le port UART2 de l'ESP32.
    \item Test des commandes AT de base : \texttt{AT}, \texttt{AT+CSQ} (force du signal), \texttt{AT+CMGF=1} (mode texte SMS).
    \item Écriture de la fonction qui lit les SMS reçus et détecte les commandes : ``\texttt{START}'', ``\texttt{STOP}'', ``\texttt{RESET}''.
    \item Test d'envoi d'un SMS depuis un téléphone pour démarrer le chronomètre à distance.
\end{itemize}

\newpage

% ==============================================================================
\section{Semaines 5 \& 6 -- Assemblage, tests et documentation (Jours 25 à 30)}
% ==============================================================================

\colorbox{semaine5}{\parbox{\linewidth}{\centering\textbf{Objectif : Monter le boîtier, tester le système complet et rédiger la documentation}}}

\bigskip

\subsection*{Jour 25 -- Préparation du châssis}
\begin{itemize}
    \item Récupération du vieux châssis de chronomètre à LED.
    \item Découpe et préparation des logements pour les 42 barres LED COB 24\,V.
    \item Organisation de la disposition des segments dans le châssis (segments a, b, c, d, e, f, g).
    \item Passage des fils à l'arrière du châssis pour garder la face avant propre.
\end{itemize}

\subsection*{Jour 26 -- Installation des rubans LED}
\begin{itemize}
    \item Découpe des rubans LED COB en barres de 12\,cm.
    \item Soudure directe des fils positif (+) et négatif (--) sur chaque barre LED.
    \item Méthode ``bouquet'' : regroupement des 7 fils rouges de chaque chiffre en un seul fil commun.
    \item Test de chaque segment individuellement avant la fermeture du boîtier.
\end{itemize}

\subsection*{Jour 27 -- Installation du diffuseur optique}
\begin{itemize}
    \item Découpe d'un panneau en PMMA opaque blanc de 3\,mm d'épaisseur.
    \item Fixation du panneau PMMA sur la face avant du boîtier.
    \item Vérification que les segments sont bien visibles à travers le diffuseur.
    \item Réglage de la distance entre les LED et le diffuseur pour éviter les points lumineux.
\end{itemize}

\subsection*{Jour 28 -- Tests électriques et mesures}
\begin{itemize}
    \item Connexion du PCB au châssis et à l'alimentation 24\,V.
    \item Mesure de la tension et du courant total avec un multimètre.
    \item Test du condensateur $C_1$ avec un oscilloscope : vérification que la tension ne chute pas de plus de 80\,mV lors des transmissions GSM.
    \item Vérification thermique des ULN2803A après 30 minutes de fonctionnement continu.
\end{itemize}

\subsection*{Jour 29 -- Tests du système complet}
\begin{itemize}
    \item Test complet du chronomètre : démarrage, arrêt, remise à zéro depuis l'interface web.
    \item Test de la commande par SMS depuis un téléphone à distance.
    \item Vérification de la lisibilité de l'affichage en pleine lumière et à plus de 20 mètres.
    \item Correction des derniers problèmes logiciels (retard d'affichage, réponse web lente).
\end{itemize}

\subsection*{Jour 30 -- Rédaction de la documentation et clôture du stage}
\begin{itemize}
    \item Rédaction du rapport de stage complet sous \LaTeX{} (avec schémas, tableaux, calculs).
    \item Mise en ligne du projet sur GitHub (\texttt{github.com/becem0}) avec tous les fichiers sources.
    \item Rédaction de ce journal de stage.
    \item Bilan personnel : compétences acquises, difficultés rencontrées, points à améliorer.
\end{itemize}

\newpage

% ==============================================================================
\section*{Bilan du stage}
\addcontentsline{toc}{section}{Bilan du stage}
% ==============================================================================

\subsection*{Compétences acquises}
\begin{itemize}
    \item Maîtrise du logiciel de conception de PCB \textbf{KiCad 10}.
    \item Calcul et dimensionnement de circuits électroniques de puissance (alimentation, filtres, transistors).
    \item Programmation d'un microcontrôleur \textbf{ESP32} en langage C++ (Arduino).
    \item Création d'un serveur web embarqué et d'une interface HTML/CSS/JavaScript.
    \item Utilisation des commandes AT pour contrôler un module GSM.
    \item Travail pratique : soudure, câblage, montage mécanique.
\end{itemize}

\subsection*{Difficultés rencontrées}
\begin{itemize}
    \item Gestion des pics de courant du module SIM800L (problème de brownout réglé avec $C_1$).
    \item Contraintes de fabrication artisanale du PCB (trous non métallisés, NPTH) qui ont nécessité une adaptation du design avec des sockets usinées.
    \item Problèmes de boot de la carte ESP32 clone (watchdog reset) -- nécessite une carte de meilleure qualité.
    \item Organisation du câblage à l'arrière du châssis avec 42 × 2 fils = 84 fils au total.
\end{itemize}

\subsection*{Conclusion personnelle}
Ce stage m'a permis de réaliser un projet complet de A à Z, depuis l'étude théorique jusqu'à la fabrication et les tests. J'ai appris à travailler de façon méthodique, à résoudre des problèmes techniques réels et à documenter mon travail de manière professionnelle. Ce projet m'a donné confiance en mes capacités à mener un projet d'ingénierie de bout en bout.

\bigskip
\begin{center}
\rule{0.4\linewidth}{0.4pt}\\[0.3cm]
{\small Becem Ameur -- ENIM, Monastir, Tunisie -- 2026}
\end{center}

\end{document}
